\chapter{Att kontrollera ett påstående}
\label{app:verifiera}

Det här materialet gör, som allt undervisningsmaterial, en rad
påståenden.  De flesta kan ni pröva själva genom att köra programmen ---
det är hela poängen med att koden på bilderna är samma kod som körs.  Men
några påståenden går inte att pröva på det sättet:
\begin{enumerate}
  \item att en uppslagning i en uppslagslista i genomsnitt kostar lika
    mycket oavsett hur många poster det finns, medan samma sökning i en
    lista kostar mer ju längre listan är (\cref{sec:valja});
  \item att en nyckel måste vara oföränderlig, och att uppslagslistan
    beter sig som materialet beskriver i övrigt
    (\cref{sec:nycklar,sec:saknad});
  \item att missuppfattningar om listreferenser förekommer hos nybörjare
    (\cref{sec:uppslagslistor});
  \item att en funktion ska ha ett ansvar och att lösningen ska hållas
    enkel --- principerna \foreignlanguage{english}{SRP} och
    \foreignlanguage{english}{KISS} (\cref{sec:valja,sec:telefonbok});
  \item de didaktiska påståenden som anteckningarna i marginalen gör om
    hur människor lär sig.
\end{enumerate}
Hur vet man om sådant stämmer?  Samma fråga möter er varje gång ni läser
en lärobok, en webbsida eller ett svar från en språkmodell --- och varje
gång ni själva ska skriva en rapport.  Den här bilagan beskriver en
arbetsgång för att kontrollera ett påstående.
\Cref{tab:dicts-pastaenden} visar var vart och ett av materialets
påståenden är belagt.  Lägg särskilt märke till det första: en siffra om
hur snabb en operation är ser exakt ut som ett faktum, men gäller bara
under förutsättningar som måste anges.

\section{Gör påståendet till en fråga}

Ett påstående går att kontrollera först när det är tydligt vad som skulle
kunna visa att det är \emph{fel}.  Formulera det därför som en fråga med
ett svar: \enquote{tar det längre tid att söka i en lista ju längre listan
är?} kan besvaras med ja eller nej, medan \enquote{listor är viktiga} inte
kan det --- ingen tänkbar källa skulle visa att det är fel, eftersom
\enquote{viktigt} är ett omdöme och inte ett faktum.  Ofta döljer sig
flera påståenden i ett: \enquote{stegvis förfining är en etablerad metod
som härrör från Wirth} --- ett påstående ni mötte i föreläsningen
\emph{Algoritmiskt tänkande} --- är två frågor, om \emph{ursprung} och om
\emph{etablering}, och de kräver olika slags källor.

\section{Välj rätt sorts källa}

Olika frågor besvaras av olika källor.
\begin{description}
  \item[Uppslagsverk] för vad ett ord \emph{betyder}.
  \item[Primärkällan] för var en idé \emph{kommer ifrån}, och för hur ett
    system \emph{beter sig}.  Frågan om vad en av Pythons behållare gör
    besvaras av Pythons egen dokumentation, inte av en lärobok som
    återger den.
  \item[Översikter] (\foreignlanguage{english}{reviews}) för om något är
    \emph{etablerat}, eller för vad forskningen \emph{sammantaget} säger.
  \item[Enskilda studier] för ett specifikt resultat --- med förbehållet
    att en enskild studie bara visar att något observerats en gång.
\end{description}
Var man hittar dem: universitetsbibliotekets söktjänst, och databaser som
Scopus, Web of Science, IEEE Xplore, ACM Digital Library, DBLP (datalogi)
och OpenAlex (öppen, och trots namnet inte detsamma som öppen tillgång).
Google Scholar täcker brett men rankar ogenomskinligt.

\section{Sök så att du kan ha fel}

Den som söker på namn hon redan känner till hittar det hon väntar sig.
Sök därför \emph{ämnesdrivet} --- på begrepp, inte författare.  Några
praktiska regler: håll frågorna korta, ställ samma frågor i flera
databaser, sök begreppets alla namn, och sök också efter
\emph{invändningar} genom att lägga till ord som \enquote{critique},
\enquote{limitations} eller \enquote{failure}.  Ett påstående som
överlevt en motsökning är värt mer än ett som aldrig prövats, och en
invändning som hittas blir ett \emph{förbehåll} i svaret, inte ett
nederlag.

\section{Läs i fulltext, och skriv ner hur}

En träff i en databas är inte ett belägg, och sammanfattningen
(\foreignlanguage{english}{abstract}) räcker inte heller: den säger vad
författarna vill ha sagt, inte vad de visat.  Läs hela texten och leta
upp \emph{det ställe} som styrker påståendet --- ett ordagrant citat med
sidnummer.  Räkna aldrig en källa ni inte läst.

Skriv sedan ner hur ni gick till väga, så att någon annan kan följa
spåret.  I det här materialet står den anteckningen som en kommentar
ovanför varje post i litteraturlistans källfil, med fasta fält:
\begin{description}
  \item[\texttt{CLAIM}] vilket påstående källan ska styrka;
  \item[\texttt{FOUND-VIA}] hur den hittades;
  \item[\texttt{PICKED}] varför just den, framför andra träffar;
  \item[\texttt{QUOTE}] det ordagranna citatet, med sida eller avsnitt;
  \item[\texttt{VERIFIED}] att fulltexten lästs, och var;
  \item[\texttt{COUNTER}] vad motsökningen gav;
  \item[\texttt{DATE}] när.
\end{description}
Det tar ett par minuter per källa, och det är det enda som gör det
möjligt att om ett år svara på frågan \enquote{hur vet du det?}

\section{Dra slutsatsen --- med förbehåll}
\label{sec:sok-slutsats}

Svaret skrivs ut tillsammans med det som talar emot.  Ett påstående om
hur snabb en operation är gäller till exempel en viss implementation och
ett genomsnittsfall, inte alla tänkbara fall --- och det förbehållet hör
till svaret.  Förbehållen är inte svagheter i svaret.  De \emph{är}
svaret, och de avgör hur påståendet får användas i materialet.

\section{En anmärkning om språkmodeller}

En språkmodell kan användas för att \emph{hitta} och \emph{sortera}
källor.  Lägg märke till arbetsfördelningen: varje källa som citeras ska
läsas och väljas av en människa.  Använd gärna språkmodeller för att
hitta och sortera; låt dem aldrig vara det sista ledet.  De kan göra
fel, men det är ni själva som måste stå till svars för innehållet.

\section{Materialets påståenden och var de är belagda}
\label{sec:dicts-pastaenden}

Inget av den här modulens påståenden krävde en egen litteratursökning.
Frågor om vad Python gör besvaras av Pythons egen
dokumentation\autocite{PythonDataStructures}, som är primärkällan.
Prestandapåståendet i \cref{sec:valja} är också en fråga om Python och
inte om forskningsläget: det besvaras av projektets egen
sammanställning av operationernas kostnad\autocite{PythonTimeComplexity},
där \enquote{x in s} anges som \(O(n)\) för en lista och som \(O(1)\) i
genomsnitt för en uppslagslista.  Alternativet --- att ta tid på
operationerna --- hade mätt den här datorn och inte datastrukturen.
Sidan säger själv att den beskriver \enquote{current CPython} och att
siffrorna är genomsnitt, och det förbehållet står därför också i texten.

Att missuppfattningar om listreferenser förekommer hos nybörjare kommer
från Veerasamy med fleras\autocite{KumarVeerasamy2016} studie av en
tentamen i introduktionsprogrammering.  Den artikeln har inte kunnat
läsas i fulltext härifrån, så materialet lutar sig bara mot
sammanfattningen och påstår ingenting om hur vanliga
missuppfattningarna är.  Principen att välja den enklaste behållare som
klarar de operationer man behöver (\foreignlanguage{english}{KISS}) och
principen att en funktion ska ha ett ansvar
(\foreignlanguage{english}{SRP}) namnges här; det vetenskapliga
underlaget finns i föreläsningen \emph{Funktioner}, bilaga~C
respektive~B.  Att kontrast måste upplevas samtidigt för att en
aspekt ska kunna urskiljas är hämtat från Martons \emph{Necessary
Conditions of Learning}\autocite[45--47]{Marton2015}.

\begin{table}[!htb]
  \centering
  \caption{Modulens påståenden, som frågor, och var svaret finns.}
  \label{tab:dicts-pastaenden}
  \begin{tabular}{@{}p{0.46\linewidth}p{0.44\linewidth}@{}}
    \toprule
    Fråga & Var det är belagt \\
    \midrule
    Kostar en uppslagning i en uppslagslista lika mycket oavsett antalet
    poster, till skillnad från en sökning i en lista?
      & Primärkälla: Pythons egen sammanställning av operationernas
        kostnad, läst i fulltext 2026-09-03.  Ja --- i genomsnitt, i
        CPython; i värsta fall nej. \\
    Beter sig uppslagslistan i övrigt som materialet beskriver, och
    måste en nyckel vara oföränderlig?
      & Primärkälla: Pythons egen dokumentation, läst i fulltext
        2026-09-03. \\
    Förekommer missuppfattningar om listreferenser hos nybörjare?
      & Ja, enligt sammanfattningen av Veerasamy med flera; fulltexten
        har inte kunnat läsas härifrån. \\
    Ska en funktion ha ett ansvar, och ska lösningen hållas enkel?
      & Redan belagt i föreläsningen \emph{Funktioner}, bilaga~B
        respektive~C. \\
    Måste kontrast upplevas samtidigt för att en aspekt ska kunna
    urskiljas?
      & Primärkälla: Marton, \emph{Necessary Conditions of Learning},
        s.~45--47. \\
    \bottomrule
  \end{tabular}
\end{table}
